% Options for packages loaded elsewhere
% Options for packages loaded elsewhere
\PassOptionsToPackage{unicode}{hyperref}
\PassOptionsToPackage{hyphens}{url}
\PassOptionsToPackage{dvipsnames,svgnames,x11names}{xcolor}
%
\documentclass[
  letterpaper,
  DIV=11,
  numbers=noendperiod]{scrartcl}
\usepackage{xcolor}
\usepackage{amsmath,amssymb}
\setcounter{secnumdepth}{-\maxdimen} % remove section numbering
\usepackage{iftex}
\ifPDFTeX
  \usepackage[T1]{fontenc}
  \usepackage[utf8]{inputenc}
  \usepackage{textcomp} % provide euro and other symbols
\else % if luatex or xetex
  \usepackage{unicode-math} % this also loads fontspec
  \defaultfontfeatures{Scale=MatchLowercase}
  \defaultfontfeatures[\rmfamily]{Ligatures=TeX,Scale=1}
\fi
\usepackage{lmodern}
\ifPDFTeX\else
  % xetex/luatex font selection
\fi
% Use upquote if available, for straight quotes in verbatim environments
\IfFileExists{upquote.sty}{\usepackage{upquote}}{}
\IfFileExists{microtype.sty}{% use microtype if available
  \usepackage[]{microtype}
  \UseMicrotypeSet[protrusion]{basicmath} % disable protrusion for tt fonts
}{}
\makeatletter
\@ifundefined{KOMAClassName}{% if non-KOMA class
  \IfFileExists{parskip.sty}{%
    \usepackage{parskip}
  }{% else
    \setlength{\parindent}{0pt}
    \setlength{\parskip}{6pt plus 2pt minus 1pt}}
}{% if KOMA class
  \KOMAoptions{parskip=half}}
\makeatother
% Make \paragraph and \subparagraph free-standing
\makeatletter
\ifx\paragraph\undefined\else
  \let\oldparagraph\paragraph
  \renewcommand{\paragraph}{
    \@ifstar
      \xxxParagraphStar
      \xxxParagraphNoStar
  }
  \newcommand{\xxxParagraphStar}[1]{\oldparagraph*{#1}\mbox{}}
  \newcommand{\xxxParagraphNoStar}[1]{\oldparagraph{#1}\mbox{}}
\fi
\ifx\subparagraph\undefined\else
  \let\oldsubparagraph\subparagraph
  \renewcommand{\subparagraph}{
    \@ifstar
      \xxxSubParagraphStar
      \xxxSubParagraphNoStar
  }
  \newcommand{\xxxSubParagraphStar}[1]{\oldsubparagraph*{#1}\mbox{}}
  \newcommand{\xxxSubParagraphNoStar}[1]{\oldsubparagraph{#1}\mbox{}}
\fi
\makeatother


\usepackage{longtable,booktabs,array}
\usepackage{calc} % for calculating minipage widths
% Correct order of tables after \paragraph or \subparagraph
\usepackage{etoolbox}
\makeatletter
\patchcmd\longtable{\par}{\if@noskipsec\mbox{}\fi\par}{}{}
\makeatother
% Allow footnotes in longtable head/foot
\IfFileExists{footnotehyper.sty}{\usepackage{footnotehyper}}{\usepackage{footnote}}
\makesavenoteenv{longtable}
\usepackage{graphicx}
\makeatletter
\newsavebox\pandoc@box
\newcommand*\pandocbounded[1]{% scales image to fit in text height/width
  \sbox\pandoc@box{#1}%
  \Gscale@div\@tempa{\textheight}{\dimexpr\ht\pandoc@box+\dp\pandoc@box\relax}%
  \Gscale@div\@tempb{\linewidth}{\wd\pandoc@box}%
  \ifdim\@tempb\p@<\@tempa\p@\let\@tempa\@tempb\fi% select the smaller of both
  \ifdim\@tempa\p@<\p@\scalebox{\@tempa}{\usebox\pandoc@box}%
  \else\usebox{\pandoc@box}%
  \fi%
}
% Set default figure placement to htbp
\def\fps@figure{htbp}
\makeatother


% definitions for citeproc citations
\NewDocumentCommand\citeproctext{}{}
\NewDocumentCommand\citeproc{mm}{%
  \begingroup\def\citeproctext{#2}\cite{#1}\endgroup}
\makeatletter
 % allow citations to break across lines
 \let\@cite@ofmt\@firstofone
 % avoid brackets around text for \cite:
 \def\@biblabel#1{}
 \def\@cite#1#2{{#1\if@tempswa , #2\fi}}
\makeatother
\newlength{\cslhangindent}
\setlength{\cslhangindent}{1.5em}
\newlength{\csllabelwidth}
\setlength{\csllabelwidth}{3em}
\newenvironment{CSLReferences}[2] % #1 hanging-indent, #2 entry-spacing
 {\begin{list}{}{%
  \setlength{\itemindent}{0pt}
  \setlength{\leftmargin}{0pt}
  \setlength{\parsep}{0pt}
  % turn on hanging indent if param 1 is 1
  \ifodd #1
   \setlength{\leftmargin}{\cslhangindent}
   \setlength{\itemindent}{-1\cslhangindent}
  \fi
  % set entry spacing
  \setlength{\itemsep}{#2\baselineskip}}}
 {\end{list}}
\usepackage{calc}
\newcommand{\CSLBlock}[1]{\hfill\break\parbox[t]{\linewidth}{\strut\ignorespaces#1\strut}}
\newcommand{\CSLLeftMargin}[1]{\parbox[t]{\csllabelwidth}{\strut#1\strut}}
\newcommand{\CSLRightInline}[1]{\parbox[t]{\linewidth - \csllabelwidth}{\strut#1\strut}}
\newcommand{\CSLIndent}[1]{\hspace{\cslhangindent}#1}



\setlength{\emergencystretch}{3em} % prevent overfull lines

\providecommand{\tightlist}{%
  \setlength{\itemsep}{0pt}\setlength{\parskip}{0pt}}



 


\KOMAoption{captions}{tableheading}
\usepackage{fancyvrb}
\usepackage{fvextra}
\RecustomVerbatimEnvironment{verbatim}{Verbatim}{breaklines=true,breakanywhere=true,fontsize=\small}
\makeatletter
\@ifpackageloaded{caption}{}{\usepackage{caption}}
\AtBeginDocument{%
\ifdefined\contentsname
  \renewcommand*\contentsname{Table of contents}
\else
  \newcommand\contentsname{Table of contents}
\fi
\ifdefined\listfigurename
  \renewcommand*\listfigurename{List of Figures}
\else
  \newcommand\listfigurename{List of Figures}
\fi
\ifdefined\listtablename
  \renewcommand*\listtablename{List of Tables}
\else
  \newcommand\listtablename{List of Tables}
\fi
\ifdefined\figurename
  \renewcommand*\figurename{Figure}
\else
  \newcommand\figurename{Figure}
\fi
\ifdefined\tablename
  \renewcommand*\tablename{Table}
\else
  \newcommand\tablename{Table}
\fi
}
\@ifpackageloaded{float}{}{\usepackage{float}}
\floatstyle{ruled}
\@ifundefined{c@chapter}{\newfloat{codelisting}{h}{lop}}{\newfloat{codelisting}{h}{lop}[chapter]}
\floatname{codelisting}{Listing}
\newcommand*\listoflistings{\listof{codelisting}{List of Listings}}
\makeatother
\makeatletter
\makeatother
\makeatletter
\@ifpackageloaded{caption}{}{\usepackage{caption}}
\@ifpackageloaded{subcaption}{}{\usepackage{subcaption}}
\makeatother
\usepackage{bookmark}
\IfFileExists{xurl.sty}{\usepackage{xurl}}{} % add URL line breaks if available
\urlstyle{same}
\hypersetup{
  pdftitle={xtitle/AI},
  pdfauthor={st. schwarz},
  colorlinks=true,
  linkcolor={blue},
  filecolor={Maroon},
  citecolor={Blue},
  urlcolor={Blue},
  pdfcreator={LaTeX via pandoc}}


\title{xtitle/AI}
\usepackage{etoolbox}
\makeatletter
\providecommand{\subtitle}[1]{% add subtitle to \maketitle
  \apptocmd{\@title}{\par {\large #1 \par}}{}{}
}
\makeatother
\subtitle{AI and language change \& variation}
\author{st. schwarz}
\date{2026-03-03}
\begin{document}
\maketitle


\section{index}\label{index}

\section{abstract B: idee 2}\label{abstract-b-idee-2}

language is developing and changing. there are many factors that
influence language change, but i want to focus on one that may become
relevant in todays language development: the change of language by
constant and increasing use of AI tools that :communicate: with us as
partners one may consider :real: or human-equivalent.

\subsection{inspiration}\label{inspiration}

now what are we experiencing if communicating with an artificial
intelligence? first comes to mind the seemingly :natural language:
adressed to us. one may feel as if talking to a human when asking
questions and getting a response. studies prove that a significant
amount of us show behaviour towards the AI that one would expect humans
show only towards each other. that leads to the first question:

\begin{quote}
if we hold the AI as a human communication partner, could its behaviour
towards us (here: their language) influence the way we talk/act
viceversa? can people learn from an AI \emph{how} to talk and what would
they learn in this case? what is the language \emph{taught} here
specifically? do we adapt to patterns or linguistic markers common for
\emph{AI speech}?
\end{quote}

\subsection{AI speech: wtf}\label{ai-speech-wtf}

we can assume as common ground that the language used (here: output) by
LLMs seems rather neutral, deprived of features deviating from the norm.
its rather easy to understand, doesnt contain irony or sarcasm very
often nor hyperbolic sentence structures (if not explicitly prompted)
and could be very well used in a textbook for learners. it may be
considered \emph{universal} in aspects of transferability into other
languages. it uses to not contain any specific vocabular or non-standard
phrases. the syntax and grammar seems to follow the corresponding rules
as the models are trained on large corpora of natural language. if we
would (and we will do that) analyse a corpus of LLM outputs we very
probably will find that in any feature it complies with the average
feature matrix of any language compared. so if one language goes like
SPO with having an average wordcount of 5wds/phrase and an average
wordlength of 5 chars then the LLM certainly will show the same features
for output in that language. no magic so far.

But: what if learners or people with deficient language skills begin to
sync their output with the artificial language in their chatverlauf?
simple like: beginning a response firstly with an appreciation of the
:very interesting question: whatever the other may have asked? we're
already heading that way\ldots{}

There may also be tiny (oberlehrerhafte) standardisations of our own
speech peculiarities (idiosyncrasies) we are confronted with which we
are kind of nudged to relativate if always sending them into a black
hole.

\subsection{methods}\label{sec-methods}

first to do would be to create (or search for) a corpus of AI generated
output. to use an existing corpus would prevent biasing which on the
other hand could be interesting to explore i.e.~we could by building
(generating) a corpus ourselves on the basis of certain dedicated
prompts\footnote{which could very well be adapted to our research
  question} force the AI to generate phenomena of interest to our
research question. where we get into medias res\ldots{}

\subsection{focusing questions}\label{focusing-questions}

\begin{enumerate}
\def\labelenumi{\arabic{enumi}.}
\tightlist
\item
  does a generative AI generally produces output that is in any aspects
  of interest for linguistic research?
\item
  how will users prime the output?
\item
  how are users adapting their own production to the output?

  \begin{enumerate}
  \def\labelenumii{\arabic{enumii}.}
  \tightlist
  \item
    is there any consistency concerning this adaptation?
  \item
    is there then societal adaptation of AI produced language?
  \item
    what are the rules (historic evidence) for adaptation?
  \end{enumerate}
\end{enumerate}

\subsection{going deeper}\label{going-deeper}

\subsubsection{corpus creation}\label{corpus-creation}

as proposed in Section~\ref{sec-methods} a secure way to building a
corpus of AI speech - which we need to explore phenomena - is to archive
LLM output to dedicated prompts. we will use an open source model
provided with llama which can run on a local machine and is adressed via
an API by script.

\paragraph{constraints}\label{constraints}

we will form prompts following these directives lead by our research
question: 1. correction of mistakes to devise actual ``knowledge'' of
model concerning standardisation and normative aspects

\section{capacities}\label{capacities}

to arrive at a research question, maybe discarding above

\subsection{what i would like to\ldots{}}\label{what-i-would-like-to}

\begin{itemize}
\tightlist
\item
  AI chat queries analysis
\item
  analyse linguistic knowledge capacities of AIs
\item
  proofread responses, factchecking
\end{itemize}

\subsection{\texorpdfstring{and what i actually \textbf{can}
do}{and what i actually can do}}\label{and-what-i-actually-can-do}

\begin{itemize}
\tightlist
\item
  statistics
\item
  automated prompting using APIs
\item
  automated response processing
\item
  masked prompt generation
\end{itemize}

\section{annotations}\label{annotations}

the research-question-still-to-develop will orientate on the findings
(\href{qa.qmd\#tbl_qa}{Table}, my annotations) in BSI
(\citeproc{ref-bsi_wie_2025}{2025}) where i see some interesting points
although the source itself may be subject to doubt.

\textbf{prompt-linguistics} seems to be a promising keyword in this
context, also Vechta U (\citeproc{ref-vechta_u_beyond_2024}{2024})
published a CfP\footnote{for which i in the moment dont find a reliable
  source anymore so its not cited.} in that genre for a themenheft which
will be published april 2026 titled: \textbf{Themenheft Beyond
Prompting?! Sozio-technische Systeme, KI und Medienbildung in der
Post-Digitalität}, edited by Annekatrin Bock, Lina Franken, Franco Rau,
Jessica Kühn und Ada Fehr.

\subsection{from here\ldots{}}\label{from-here}

\ldots we come to a maybe more feasable to explore topic within the
research of semantic, syntactic and pragmatic features of speech of AI
users that is affected by that use. in the not citable BSI
(\citeproc{ref-bsi_wie_2025}{2025}) a focus lies on how users language
changes depending on their making ``heavy use'' of AI tools. there seems
to be an influence of the affordances to create successful/optimized
prompts for best results on a. the language used in that prompts and
furthermore b. the language beyond that usecase say the users everyday
life. which is exactly what we were looking for, and from here, we can
dig into corpora to find manifestations of features the study discovered
as e.g.~simplification of complex syntactic structures in favour of
reference based patterns.

\subsubsection{Q}\label{q}

\paragraph{\texorpdfstring{characteristics of the optimal prompt,
Leidinger, Rooij, and Shutova
(\citeproc{ref-leidinger_language_2023}{2023})}{characteristics of the optimal prompt, Leidinger, Rooij, and Shutova (2023)}}\label{characteristics-of-the-optimal-prompt-leidinger_language_2023}

\begin{itemize}
\tightlist
\item
  can we trace these in corpora?
\item
  findings:

  \begin{itemize}
  \tightlist
  \item
    no significant accuracy increase with simplification or lower
    perplexity score
  \item
    no accuracy increase with more frequent synonymes
  \end{itemize}
\item
  ie: there doesnt seem to be a generally ``optimal language'' to
  prompting

  \begin{itemize}
  \tightlist
  \item
    may not seek for structural features really
  \end{itemize}
\item
  contradictory to BSI study and intuitive guess (experience) on
  optimized prompt language
\end{itemize}

\paragraph{\texorpdfstring{Empirical evidence of Large Language Model's
influence on human spoken communication, Yakura et al.
(\citeproc{ref-yakura_empirical_2025}{2025})}{Empirical evidence of Large Language Model's influence on human spoken communication, Yakura et al. (2025)}}\label{empirical-evidence-of-large-language-models-influence-on-human-spoken-communication-yakura_empirical_2025}

\begin{itemize}
\tightlist
\item
  finds significant changes in human spoken communication past GPT onset
\item
  try replication on german language:

  \begin{itemize}
  \tightlist
  \item
    \href{yakura.html}{preview motivation and preliminary}
  \item
    \href{gemini.html}{preview paper \& first results}
  \end{itemize}
\end{itemize}

\section*{References}\label{references}
\addcontentsline{toc}{section}{References}

\phantomsection\label{refs}
\begin{CSLReferences}{1}{0}
\bibitem[\citeproctext]{ref-abrami_german_2022}
Abrami, Giuseppe, Mevlüt Bagci, Leon Hammerla, and Alexander Mehler.
2022. {``German {Parliamentary} {Corpus} ({GerParCor}).''} In
\emph{Proceedings of the {Thirteenth} {Language} {Resources} and
{Evaluation} {Conference}}, edited by Nicoletta Calzolari, Frédéric
Béchet, Philippe Blache, Khalid Choukri, Christopher Cieri, Thierry
Declerck, Sara Goggi, et al., 1900--1906. Marseille, France: European
Language Resources Association.
\url{https://aclanthology.org/2022.lrec-1.202/}.

\bibitem[\citeproctext]{ref-aleksic_opinion_2025}
Aleksic, Adam. 2025. {``Opinion {\textbar} {It}'s Happening: {People}
Are Starting to Talk Like {ChatGPT}.''} \emph{The Washington Post},
August.
\url{https://www.washingtonpost.com/opinions/2025/08/20/chatgpt-claude-chatbots-language/}.

\bibitem[\citeproctext]{ref-noauthor_automatic_nodate}
{``Automatic Register Identification for the Open Web Using Multilingual
Deep Learning.''} n.d. Accessed January 3, 2026.
\url{https://arxiv.org/html/2406.19892v3}.

\bibitem[\citeproctext]{ref-bates_fitting_2015}
Bates, Douglas, Martin Mächler, Ben Bolker, and Steve Walker. 2015.
{``Fitting {Linear} {Mixed}-{Effects} {Models} {Using} Lme4.''}
\emph{Journal of Statistical Software} 67 (1): 1--48.
\url{https://doi.org/10.18637/jss.v067.i01}.

\bibitem[\citeproctext]{ref-bsi_wie_2025}
BSI, Brand Science Institute. 2025. {``Wie {KI} Unsere {Sprache}
Verändert -- {Eine} Empirische {Studie}.''}
\url{https://www.bsi.ag/cases/104-case-studie-wie-ki-unsere-sprache-veraendert---eine-empirische-studie.html}.

\bibitem[\citeproctext]{ref-dip_dip_2026}
DIP. 2026. {``{DIP} - {Bundestagsprotokolle}.''} Docs. \emph{DIP - API}.
Berlin.
\url{https://dip.bundestag.de/\%C3\%BCber-dip/hilfe/api\#content}.

\bibitem[\citeproctext]{ref-duke_upress_critical_2025}
Duke UPress. 2025. {``Critical {AI}.''}
\url{https://www.dukeupress.edu/critical-ai}.

\bibitem[\citeproctext]{ref-flach_collostructions_2021}
Flach, Susanne. 2021. {``Collostructions: {An} {R} Implementation for
the Family of Collostructional Methods. {Package} Version v.0.2.0.''}
\emph{GitHub - Skeptikantin}. \url{https://github.com/skeptikantin}.

\bibitem[\citeproctext]{ref-fobbe_r_2026}
Fobbe, Sean. 2026. {``{[}{R}{]} {Source} {Code} Des '{Corpus} Der
{Plenarprotokolle} Des {Deutschen} {Bundestages}'
({CPP}-{BT}-{Source}).''} Zenodo.
\url{https://doi.org/10.5281/zenodo.18177197}.

\bibitem[\citeproctext]{ref-kalwa_noch_2025}
Kalwa, Nina. 2025. {``»{Noch} Steckt {KI} in Den {Kinderschuhen}«.''}
\emph{Zeitschrift Für Literaturwissenschaft Und Linguistik} 55 (2):
379--405. \url{https://doi.org/10.1007/s41244-025-00379-0}.

\bibitem[\citeproctext]{ref-kraus_studie_2025}
Krauß, Patrick, and FAU U. 2025. {``Studie Zeigt: {KI} Lernt
{Sprachregeln} Beim {Lesen}.''} \emph{FAU}.
\url{https://www.fau.de/2025/11/news/forschung/studie-ki-lernt-sprachregeln-beim-lesen/}.

\bibitem[\citeproctext]{ref-leidinger_language_2023}
Leidinger, Alina, Robert van Rooij, and Ekaterina Shutova. 2023. {``The
Language of Prompting: {What} Linguistic Properties Make a Prompt
Successful?''} In \emph{Findings of the {Association} for
{Computational} {Linguistics}: {EMNLP} 2023}, edited by Houda Bouamor,
Juan Pino, and Kalika Bali, 9210--32. Singapore: Association for
Computational Linguistics.
\url{https://doi.org/10.18653/v1/2023.findings-emnlp.618}.

\bibitem[\citeproctext]{ref-milicka_ai_2025}
Milička, Jiří, Anna Marklová, and Václav Cvrček. 2025. {``{AI} {Brown}
and {AI} {Koditex}: {LLM}-{Generated} {Corpora} {Comparable} to
{Traditional} {Corpora} of {English} and {Czech} {Texts}.''}
\url{https://doi.org/10.48550/arxiv.2509.22996}.

\bibitem[\citeproctext]{ref-ramirez_chatgpt_2025}
Ramirez, Vanessa Bates. 2025. {``{ChatGPT} {Is} {Changing} the {Words}
{We} {Use} in {Conversation}.''} \emph{Scientific American}.
\url{https://www.scientificamerican.com/article/chatgpt-is-changing-the-words-we-use-in-conversation/}.

\bibitem[\citeproctext]{ref-ror_ror_2025}
ROR, Research Organization Registry. 2025. {``{ROR} {Data}.''} Zenodo.
\url{https://doi.org/10.5281/zenodo.17953395}.

\bibitem[\citeproctext]{ref-roussel_einsatz_2024}
Roussel, Stephanie, and Ayaal Herdam. 2024. {``Einsatz von Künstlicher
{Intelligenz} Im Überarbeitungsprozess von {Texten} in Der
{Fremdsprache} {Deutsch}: {Fokus} Auf {Wortschatz} Und {Syntax}.''}
\emph{German as a Foreign Language} 2024 (3): 61--86.
\url{https://hal.science/hal-04999057}.

\bibitem[\citeproctext]{ref-schwarz_this_2026}
Schwarz, St. 2026. {``This Papers Evaluation Script.''}
\emph{GitHub/Esteeschwarz}. Berlin.
\url{https://github.com/esteeschwarz/SPUND-LX/blob/main/germanic/HA/LLM-003.R}.

\bibitem[\citeproctext]{ref-shrivastava_repository-level_2023}
Shrivastava, Disha, Hugo Larochelle, and Daniel Tarlow. 2023.
{``Repository-Level Prompt Generation for Large Language Models of
Code.''} In \emph{Proceedings of the 40th {International} {Conference}
on {Machine} {Learning}}, 202:31693--715. {ICML}'23. Honolulu, Hawaii,
USA: JMLR.org.

\bibitem[\citeproctext]{ref-vechta_u_beyond_2024}
Vechta U. 2024. {``Beyond {Prompting}.''}
\url{https://www.uni-vechta.de/beyondprompting/news/details/projekt-beyond-prompting-erfolgreiche-drittmitteleinwerbung}.

\bibitem[\citeproctext]{ref-wikipedia_google_2026}
Wikipedia, and Google. 2026. {``Google {Gemini}.''} \emph{Wikipedia}.
\url{https://de.wikipedia.org/w/index.php?title=Google_Gemini&oldid=263426206}.

\bibitem[\citeproctext]{ref-wu_corpus-based_2025}
Wu, JinLiang. 2025. {``A {Corpus}-{Based} {Multidimensional} {Analysis}
of {Linguistic} {Features} Between {Human}-{Authored} and
{ChatGPT}-{Generated} {Compositions}.''} \emph{International Journal of
Linguistics, Literature and Translation} 8 (5): 102--10.
\url{https://doi.org/10.32996/ijllt.2025.8.5.10}.

\bibitem[\citeproctext]{ref-yakura_empirical_2025}
Yakura, Hiromu, Ezequiel Lopez-Lopez, Levin Brinkmann, Ignacio Serna,
Prateek Gupta, Ivan Soraperra, and Iyad Rahwan. 2025. {``Empirical
Evidence of {Large} {Language} {Model}'s Influence on Human Spoken
Communication.''} arXiv.
\url{https://doi.org/10.48550/arXiv.2409.01754}.

\end{CSLReferences}




\end{document}
